\documentclass{article}

\usepackage{graphicx}
\usepackage{tikz}
\usepackage{lipsum}  % For generating dummy text
\usepackage{amsthm, amssymb, amsmath}  % For mathematical symbols

\title{Insiemi Numerici (Numeric Sets)}
\author{Davide Pulvirenti}

\begin{document}

\maketitle

\tableofcontents

\newpage

\section{Introduction}
Insiemi numerici, or numeric sets, are fundamental concepts in mathematics that play a crucial role in various branches of mathematics and its applications. A numeric set is a collection of distinct numbers or elements, and it serves as a building block for various mathematical structures and operations.

\section{Basic Definitions}
\subsection{Set Notation}
A numeric set is typically denoted by a capital letter. For example, \(A\), \(B\), or \(C\) can represent numeric sets. The elements in a set are enclosed in curly braces, like \(\{1, 2, 3, 4, 5\}\).

\subsection{Cardinality}
The cardinality of a set is the number of elements it contains. For a set \(A\), the cardinality is denoted as \(|A|\).

\subsection{Subsets}
If every element of set \(A\) is also an element of set \(B\), we say that \(A\) is a subset of \(B\), denoted as \(A \subseteq B\).

\subsection{Union and Intersection}
The union of two sets \(A\) and \(B\), denoted \(A \cup B\), is the set of all elements that belong to \(A\), \(B\), or both. The intersection of sets \(A\) and \(B\), denoted \(A \cap B\), contains elements that are in both \(A\) and \(B\).

\section{Types of Numeric Sets}
\subsection{Natural Numbers (\(\mathbb{N}\))}
The set of natural numbers includes all positive integers starting from 1, i.e., \(\mathbb{N} = \{1, 2, 3, 4, \ldots\}\).

\subsection{Integers (\(\mathbb{Z}\))}
The set of integers includes all positive and negative whole numbers, along with zero, i.e., \(\mathbb{Z} = \{\ldots, -3, -2, -1, 0, 1, 2, 3, \ldots\}\).

\subsection{Rational Numbers (\(\mathbb{Q}\))}
The set of rational numbers consists of numbers that can be expressed as fractions, where the numerator and denominator are integers, and the denominator is not zero. For example, \(\frac{3}{4}\) is a rational number.

\subsection{Real Numbers (\(\mathbb{R}\))}
The set of real numbers includes all rational and irrational numbers. It forms a continuous number line, encompassing all possible values.

\section{Conclusion}
Insiemi numerici are a fundamental concept in mathematics, providing the foundation for various mathematical operations and structures. Understanding the properties and relationships of numeric sets is crucial for many mathematical and scientific applications.

\newpage

\section{Diagram}
\begin{figure}[h]
  \centering
  \begin{tikzpicture}
    \draw (0,0) -- (2,0) -- (1,2) -- cycle;
  \end{tikzpicture}
  \caption{TikZ Diagram Example}
\end{figure}

\section{Formula}
Euler's formula, which relates the exponential function, cosine, and sine:
\begin{equation}
e^{ix} = \cos(x) + i\sin(x)
\end{equation}

\section{Table}
\begin{table}[h]
  \centering
  \begin{tabular}{|c|c|}
    \hline
    Set & Cardinality \\
    \hline
    \(\mathbb{N}\) & Countably Infinite \\
    \(\mathbb{Z}\) & Countably Infinite \\
    \(\mathbb{Q}\) & Countably Infinite \\
    \(\mathbb{R}\) & Uncountably Infinite \\
    \hline
  \end{tabular}
  \caption{Cardinalities of Numeric Sets}
\end{table}

\newpage

\begin{thebibliography}{9}
  \bibitem{set-theory-book}
  Myself,
  \textit{Introduction to Set Theory},
  Myself, Today.

  \bibitem{math-analysis-book}
  Myself,
  \textit{Mathematical Analysis: Concepts and Applications},
  Myself, Today.
\end{thebibliography}

\end{document}
